% ED26-P18-erotetic-0001: Wisniewski's propositional PMCE calculus.
\chapter{Question Evocation and the PMCE Calculus}
\label{sec:p18-question-evocation}

Premises can warrant a question while leaving its answer open. If the
available information entails that a report concerns either a defect in
the design or a defect in the implementation, but entails neither answer
separately, it warrants asking which of these descriptions applies. That
is a different achievement from establishing an answer.

This chapter gives Wi\'sniewski's calculus PMCE for that relation in
classical propositional logic \cite{Wisniewski2016QuestionEvocation}.
Its input is a finite set of declarative premises and an ordered finite
list of distinct proposed answers. Its guarantee has two parts: the
premises cover the answer list, and none of the individual answers is
already entailed. The alternatives need not be mutually exclusive.
The calculus does not claim that every evoked question is useful, cheap
to investigate, or sufficient for a larger inquiry.

\section{Declaratives, questions and evocation}

\begin{hbdefinition}{def:p18-qev-language}{Propositional language and questions}
Let \(\mathrm{Var}\) be a countably infinite set of propositional
variables. The set \(\mathrm{Form}\) is generated from these variables
by \(\neg,\lor,\land,\to,\leftrightarrow\), with parentheses fixing
the construction of a formula. There are no primitive truth constants.
A valuation \(v\) assigns \(0\) or \(1\) to each variable and extends by
the classical truth functions: negation reverses the value; conjunction
is true exactly when both arguments are true; disjunction when at least
one is true; implication is false exactly when its antecedent is true
and consequent false; equivalence is true exactly when its arguments
have the same value. In this source presentation, the same symbol
\(v\) also denotes that classical extension to formulas.

A \emph{question} is the expression
\[
 Q=?\{A_1,\ldots,A_n\},\qquad n>1,
\]
where the \(A_i\) are pairwise syntactically distinct members of
\(\mathrm{Form}\). The \emph{direct-answer set} is
\(dQ=\{A_1,\ldots,A_n\}\). A question is not itself a declarative
formula. Its displayed order is part of the expression: for example,
\(?\{p,\neg p\}\) and \(?\{\neg p,p\}\) are different questions
with the same direct-answer set. Syntactic distinctness does not require
classical inequivalence.
\index{question!direct answer}\index{erotetic logic!question syntax}
\end{hbdefinition}

\begin{hbdefinition}{def:p18-qev-entailment}{Multiple-conclusion entailment and proper entailment}
For sets \(X,Y\subseteq\mathrm{Form}\), write \(v\models X\)
when every member of \(X\) is true under \(v\). Ordinary consequence
\(X\models A\) means that every valuation satisfying \(X\) makes
\(A\) true. \emph{Multiple-conclusion entailment}, written
\(X\mathrel{\|=}Y\), means
\[
 \text{for every }v,\quad
 v\models X\ \Longrightarrow\ (\exists A\in Y)\ v(A)=1.
\]
For nonempty \(Y\), \emph{proper multiple-conclusion entailment},
written \(X\mathrel{\|\!\triangleleft}Y\), means that
\(X\mathrel{\|=}Y\) and \(X\not\models A\) for every \(A\in Y\).
These are relations on sets of declaratives; their double-bar signs have
this local entailment role, not a forcing interpretation.

The premises \(X\) \emph{evoke} the question \(Q\), written
\(E(X,Q)\), exactly when \(X\mathrel{\|\!\triangleleft}dQ\).
The first condition is \emph{coverage}; the second is
\emph{unsettledness}. Each failure of entailment has its own witnessing
valuation. One valuation need not falsify all direct answers at once.
\index{multiple-conclusion entailment}\index{proper multiple-conclusion entailment}
\index{question evocation}
\end{hbdefinition}

\begin{hbtheorem}{thm:p18-qev-semantic-profile}{proved=annex:p18.qev.semantic-profile}{Coverage, unsettledness and boundary cases}
For \(Q=?\{A_1,\ldots,A_n\}\),
\[
 E(X,Q)\quad\Longleftrightarrow\quad
 X\models A_1\lor\cdots\lor A_n
 \quad\hbox{and}\quad
 X\not\models A_i\text{ for every }i.
\]
An inconsistent \(X\) evokes no question. A question containing a
tautological answer is never evoked. Pointwise replacement of answers
by equivalent formulas preserves evocation whenever the replacement
list still forms a question. Permuting the answer list preserves
evocation, although it can change the question expression.
\end{hbtheorem}

\begin{hbexample}{def:p18-qev-settled-example}{Coverage is not enough}
Let \(p,q\) be different variables and put \(X=\{p\lor q\}\).
Every \(X\)-valuation makes \(p\) or \(q\) true. The valuations
\((p,q)=(0,1)\) and \((1,0)\) respectively witness
\(X\not\models p\) and \(X\not\models q\). Thus
\(E(X,?\{p,q\})\).

After adding \(p\), the same answer list remains exhaustive, but
\(X\cup\{p\}\models p\). It no longer evokes \(?\{p,q\}\).
By contrast, with no premises the assignment \((0,0)\) defeats
coverage. This example separates three states: uncovered alternatives,
covered but unsettled alternatives, and an already settled answer.
\end{hbexample}

\section{The exact calculus}

\begin{hbdefinition}{def:p18-qev-sequents}{Erotetic sequents and PMCE proofs}
An \emph{erotetic sequent} is \(X\vdash Q\), where \(X\) is a
finite, possibly empty set of declaratives and \(Q\) is a question.
Write \(X,B\) for \(X\cup\{B\}\), and omit \(X\) when it is
empty. The turnstile inside an erotetic sequent is a syntactic separator.
To say that the sequent is \emph{provable in PMCE} means that it labels
the root of a finite tree whose leaves are instances of the axiom schema
and whose internal nodes follow the four primitive rules below.

A \emph{literal} is a variable or its negation. Complementary literals
are \(p\) and \(\neg p\). A \emph{clause} is a literal or a
nonempty disjunction of literals; repetitions are allowed. In the axiom
schema each direct answer is a clause without a complementary pair,
whereas the disjunction of all answers contains a complementary pair.

In the sheet, \(\vec A,\vec L,\vec R\) denote displayed ordered
answer lists or list fragments. Expressions such as \(?\{\vec A,B\}\) append the
last answer to that list. All premises and conclusions of every rule
must be erotetic sequents. In particular a displayed answer list always
has at least two pairwise syntactically distinct members. The symbols
\(=\) and \(\ne\) in the syntactic provisos compare expressions,
not truth values or logical equivalence. For \(R_1,R_2\), the common
list has at least one member.
\index{PMCE}\index{erotetic sequent}
\end{hbdefinition}

\begin{hbdefinition}{def:p18-qev-rule-standing}{Derived and admissible rules}
A rule is \emph{derived} here when a fixed finite arrangement of
primitive rules transforms derivations of its premises into a derivation
of its conclusion. A rule is \emph{admissible} when provability of all
its premises guarantees provability of its conclusion. A derived rule
is admissible, but a semantic proof of admissibility does not itself
give a fixed local primitive-rule derivation.

For the schematic entries below, \(?[A\mid n]\) abbreviates
\(?\{A_1,\ldots,A_n\}\). Appended formulas are appended to that
ordered list; \(?[B\to A\mid n]\), \(?[A\mid n\to B]\) and
\(?[\neg A\mid n]\) apply the indicated operation to every answer.
Unless a rule explicitly uses the single-common-answer variant,
\(n>1\). All resulting expressions must remain questions.
\end{hbdefinition}

\begin{hbdefinition}{def:p18-qev-alpha-beta}{Conjunctive and disjunctive formula types}
The following table specifies the \(\alpha\)- and \(\beta\)-forms
used by the admissible rules. In every row,
\(\alpha\leftrightarrow(\alpha_1\land\alpha_2)\) and
\(\beta\leftrightarrow(\beta_1\lor\beta_2)\) are tautologies.
Double negation is not an \(\alpha\)-form in this classification;
the table does not purport to classify every nonliteral formula.
\begin{center}
\begin{tabular}{lll|lll}
\(\alpha\)&\(\alpha_1\)&\(\alpha_2\)&
\(\beta\)&\(\beta_1\)&\(\beta_2\)\\
\hline
\(B\land C\)&\(B\)&\(C\)&\(\neg(B\land C)\)&\(\neg B\)&\(\neg C\)\\
\(\neg(B\lor C)\)&\(\neg B\)&\(\neg C\)&\(B\lor C\)&\(B\)&\(C\)\\
\(\neg(B\to C)\)&\(B\)&\(\neg C\)&\(B\to C\)&\(\neg B\)&\(C\)
\end{tabular}
\end{center}
\end{hbdefinition}

\showcalculus{p18-question-evocation}

The primitive part of the sheet is Ax and \(R_1\)--\(R_4\).
The remaining entries are explicitly marked as derived or admissible.
\(R_4\) handles the difference between an ordered question expression
and its answer set. The rule \(R_2\) uses classical tautological
equivalence as a side condition. PMCE therefore is not an analytic
search calculus obtained just by decomposing the final question.

\begin{hbtheorem}{thm:p18-qev-soundness}{proved=annex:p18.qev.soundness}{Soundness of PMCE}
Every axiom evokes its question from the empty set. Each of the four
primitive rules preserves evocation from its premises to its conclusion.
Consequently, if \(X\vdash Q\) is provable in PMCE, then \(E(X,Q)\).
\end{hbtheorem}

\section{Normal forms and finite completeness}

\begin{hbdefinition}{def:p18-qev-normal-form}{CNF representatives and rank}
A formula is in \emph{conjunctive normal form} (CNF) when it is a
conjunction of one or more clauses. A single clause counts as a CNF
with one component. Parentheses and the selected listing of its clauses
are retained when counting components. An empty-antecedent sequent is
in normal form when every direct answer is in CNF. Its literal CNF
rank is the sum of the numbers of conjunction signs in those answers.

For a question \(?\{A_1,\ldots,A_n\}\), also consider a chosen
CNF representative \(C_i\) classically equivalent to each \(A_i\).
If \(C_i\) is displayed as a conjunction of \(m_i\geq1\) clauses,
the \emph{representative rank} of this choice is
\[
 k=\sum_{i=1}^n(m_i-1).
\]
This rank concerns the chosen representatives. It need not equal the
number of conjunction signs in the actual answer expressions, which
may contain additional tautological markers used to preserve syntactic
distinctness.
\index{question evocation!normal form}
\end{hbdefinition}

\begin{hblemma}{thm:p18-qev-classical-tools}{proved=annex:p18.qev.classical-tools}{Classical tools and syntactically distinct representatives}
Every declarative has a CNF representative. A clause is tautological
exactly when it contains complementary literals. Given a clause \(D\)
and finitely many forbidden expressions, there is a syntactically
different clause equivalent to \(D\) outside that finite collection.

Every question \(?\{A_1,\ldots,A_n\}\) has CNF answers
\(C_1,\ldots,C_n\) such that each \(C_i\) is equivalent to
\(A_i\), the \(C_i\) are pairwise syntactically distinct, and no
\(C_i\) equals any \(A_j\). Evocation is preserved by this
replacement. A derivation of the resulting question can be changed,
using only \(R_2,R_4\), into a derivation of the original question
with the same antecedent.
\end{hblemma}

\begin{hblemma}{thm:p18-qev-representative-induction}{proved=annex:p18.qev.representative-induction}{Completeness by representative rank}
Let \(Q=?\{A_1,\ldots,A_n\}\) have chosen CNF representatives
of finite representative rank \(k\). If \(E(\varnothing,Q)\),
then \(\vdash Q\) is provable in PMCE. The construction terminates
by induction on \(k\), including both cases in which the residual
conjunction after selecting a nontautological clause is, or is not,
a tautology.
\end{hblemma}

\begin{hbcorollary}{thm:p18-qev-empty-completeness}{proved=annex:p18.qev.empty-completeness}{Normal-form and empty-antecedent completeness}
Every evoked question in CNF normal form is provable from the empty
antecedent. More generally, \(E(\varnothing,Q)\) implies that
\(\vdash Q\) is provable, without a normal-form requirement on
the actual answer expressions.
\end{hbcorollary}

\begin{hblemma}{thm:p18-qev-conditionalization}{proved=annex:p18.qev.conditionalization}{Moving finite premises into the answers}
For every finite \(X\), formula \(B\), and question
\(?\{A_1,\ldots,A_n\}\),
\[
 E(X\cup\{B\},?\{A_1,\ldots,A_n\})
 \quad\Longleftrightarrow\quad
 E(X,?\{B\to A_1,\ldots,B\to A_n\}).
\]
If \(X=\{B_1,\ldots,B_m\}\), define
\(I_X(A)=B_1\to(B_2\to\cdots\to(B_m\to A)\cdots)\),
with \(I_{\varnothing}(A)=A\). Then
\[
 E(X,?\{A_1,\ldots,A_n\})
 \quad\Longleftrightarrow\quad
 E(\varnothing,?\{I_X(A_1),\ldots,I_X(A_n)\}).
\]
These transformed lists are questions: adjoining a fixed implication
prefix is injective on formula expressions.
\end{hblemma}

\begin{hbtheorem}{thm:p18-qev-completeness}{proved=annex:p18.qev.completeness}{Finite-premise soundness and completeness}
For every finite set \(X\) of declaratives and every question \(Q\),
\[
 X\vdash Q\text{ is provable in PMCE}
 \quad\Longleftrightarrow\quad E(X,Q).
\]
The completeness direction constructs a finite proof: normalize and
prove the empty-antecedent conditionalized question, then use \(R_3\)
to return each premise to the antecedent. No unrestricted left
weakening rule is used.
\end{hbtheorem}

\section{Derived and admissible rules}

\begin{hblemma}{thm:p18-qev-derived-rules}{proved=annex:p18.qev.derived-rules}{Rules at arbitrary answer positions}
The sheet's position-sensitive rules \(R_1^*,R_2^*\) and answer
transposition \(R_4^*\) are derived. Conjunction introduction or
equivalent replacement can therefore act at any selected answer position,
subject to the displayed syntactic-distinctness conditions.
\end{hblemma}

\begin{hbtheorem}{thm:p18-qev-admissible-rules}{proved=annex:p18.qev.admissible-rules}{Semantic transformations and admissibility}
All eleven semantic transformations corresponding to the source's
Facts~15--25 are admissible as shown in the sheet: right and left
complementary case removal; antecedent equivalence; reverse \(R_3\);
right-implication transformation; left \(\beta\)-combination;
both directions of left \(\alpha\)-expansion; right
\(\beta\)-expansion; right \(\alpha\)-combination; and export of
the negation of a premise as an answer.

For right \(\beta\)-expansion and right \(\alpha\)-combination,
the common answer list may consist of a single formula, provided every
displayed premise and conclusion is a question. This includes the two
small-question variants stated at the end of the source's Section~5.
The claim is admissibility for finite antecedents in PMCE, not a new
first-order calculus.
\end{hbtheorem}

\section{Examples, decision and limits}

\begin{hbtheorem}{thm:p18-qev-source-examples}{proved=annex:p18.qev.source-examples}{Five explicit PMCE derivations}
For distinct variables \(p,q,r\), each of the following is provable:
\[
\begin{array}{rcl}
 \{p\lor\neg p\}&\vdash&?\{p,\neg p\},\\
 \{p\lor q\}&\vdash&?\{p,q\},\\
 \{p\to(q\lor r),p\}&\vdash&?\{q,r\},\\
 \{(p\land q)\to r,\neg r\}&\vdash&?\{\neg p,\neg q\},\\
 \varnothing&\vdash&?\{p,q,\neg(p\lor q)\}.
\end{array}
\]
The annex supplies the primitive-rule derivations, including the
tautological equivalences needed by each \(R_2\) application.
\end{hbtheorem}

\begin{hbmethod}{def:p18-qev-decision-method}{Deciding whether the premises evoke a question}
Check first that the answer list has at least two syntactically distinct
formulas and that the antecedent is finite. List the finitely many
variables occurring in the input. Enumerate their Boolean assignments.
For each assignment satisfying the premises, check whether at least
one answer is true, and record each answer that is false at that
assignment. Return yes exactly when coverage never fails and every
answer has been recorded as false at least once. If no assignment
satisfies the premises, return no.
\end{hbmethod}

\begin{hbtheorem}{thm:p18-qev-decision}{proved=annex:p18.qev.decision}{Decision and failure of unrestricted weakening}
The method terminates after examining \(2^r\) assignments when
\(r\) variables occur in the input. It returns yes exactly when
\(E(X,Q)\), and hence exactly when \(X\vdash Q\) is provable
in PMCE. The bound counts assignments, not constant-time formula
evaluation. In general, adding premises can destroy evocation and
PMCE provability, even though it cannot destroy ordinary classical
entailment.
\end{hbtheorem}

\begin{systementry}{entry:p18-question-evocation}{PMCE: propositional question evocation}
\entryaliases{Proper multiple-conclusion entailment with erotetic sequents.}
\entrylanguage{Finite sets of classical propositional premises and ordered
questions with at least two pairwise syntactically distinct direct answers.}
\entrysemantics{Coverage of the answer set together with failure to entail
any particular answer: \(E(X,Q)\).}
\entrycalculus{sheet:p18-question-evocation}
\entryadequacy{proved=annex:p18.qev.completeness}{Sound and complete for
evocation by finite propositional antecedents.}
\entrydecidability{proved=annex:p18.qev.decision}{Finite truth-table decision
by coverage and separate countervaluations for the answers.}
\entrycharacteristic{Question order is syntactic; \(R_4\) handles
permutations. Unrestricted left weakening fails. Syntactic identity,
logical equivalence and equality of answer sets remain distinct.}
\entrytranslations{Finite antecedents can be internalized as common
conditional prefixes. Evocation is not resolution or general erotetic
implication.}
\entrysources{Wisniewski2016QuestionEvocation}
\end{systementry}

\section{Evocation, resolution and inquiry}

Evocation asks whether the premises support raising a question under the
coverage-and-unsettledness criterion. Resolution asks whether information
settles a suitable answer. In the example above, adding \(p\) advances
resolution while destroying evocation of \(?\{p,q\}\). Thus the two
judgments must not be identified.

An inference from a question, possibly together with declarative premises,
to another question is a further relation: erotetic implication, developed
in Section~\ref{sec:p18-question-to-question-implication}. PMCE's
antecedents contain only declaratives, so its completeness theorem does
not axiomatize that general relation. The source also distinguishes
evocation from a presupposition-based permission to ask in an interrogative
game. Whether a practical inquiry should ask a question depends on
additional information about goals, access to answers and costs.

The result here remains propositional. First-order questions with
quantified answer patterns require separately defined question languages
and separate completeness results. Likewise, the source's discussion of
other theories of questions, inquisitive semantics and Socratic proof
search establishes their relationship as further subjects, not an
identity of their formal systems. The truth-table decision method and
the constructive completeness proof give finite procedures; they do
not make PMCE's equivalence rule an efficient analytic search strategy.

\section{Historical position of the calculus}

Wi\'sniewski places PMCE within inferential erotetic logic, which studies
inferences with questions as conclusions, and distinguishes that program
from earlier work centered on question structure and answerhood.
In the historical account given in the source, Ajdukiewicz's 1926
proposal treats an interrogative as a sentential function closed by
an interrogative operator that restricts the variable's range. Even
polar questions receive an operator-based analysis. Kubi\'nski's later
systems investigate relations among questions, such as equivalence,
containment and relative strength, together with relations between
questions and possible answers \cite[sec.~1.3]{Wisniewski2016QuestionEvocation}.

The immediately preceding proper-multiple-conclusion calculus PMC uses
sets of declaratives on both sides of a sequent. PMCE retains a separate
question category on the right. That difference explains the explicit
structural permutation rule and the need to preserve well-formed
question expressions throughout the completeness construction. It also
explains why equality of direct-answer sets suffices for a semantic
comparison while remaining weaker than identity of question expressions.
